\section{Background}
\label{sec:background}

\textbf{Decision models.}
A decision model maps a textual \emph{state} and a typed question to a structured decision, with a probability per option.
In this work, we focus on choice questions: each specifies a set of named options, a short description of each option, and instructions for selecting among them.
Unlike a conventional language-model call that returns a generated textual response, a decision-model call returns a choice in a predefined format.
This interface can avoid generating a free-form output, but it does not, by itself, guarantee lower end-to-end latency, cost, or energy: these depend on the model, serving configuration, and downstream computation.

\jev~\cite{typesafe2026jev} is a proprietary decision model accessible through an API; its weights, architecture, and training data are not publicly available.
We evaluate version 1.13 and treat vendor claims about its performance as hypotheses to test empirically.
\kev~\cite{palmer2026kev} is a third-party, Apache-2.0-licensed implementation of the decision-model interface, available in 0.8B, 4B, and 9B parameter sizes and accepting the same API request format as \jev.
It does not use \jev's weights, and we do not assume that the two models are behaviorally equivalent.
We include \kev\ for two complementary reasons: it provides an open model that can be evaluated without access to the proprietary API, and it enables local measurements of inference latency and GPU energy that are not directly observable through a hosted API.
We therefore use \jev\ to evaluate the proprietary service and \kev\ to study the behavior and serving costs of an open, locally deployed alternative.

\textbf{Three decision points.}
We study three places where a standalone decision model could sit in an LLM system: agent control, semantic-cache guarding and model routing (Table~\ref{tab:design}).
They differ in whether the decision can remove a generation.
In agent control it cannot: the decision picks a tool or a reply, but some model must still write the tool's arguments or the reply, so a separate decider adds a call.
A cache guard's decision to reuse skips the generation, and a router's decision to stay with the small model skips the large one's.
At those two points the decider must still beat a cheaper alternative, a similarity threshold or a trained classifier, to justify its own call.

\begin{table*}[t]
\caption{The three decision points studied in this work. ``Replaces generation?'' indicates whether a decision can avoid a subsequent generative-model call. ``Cheap alternative'' identifies the baseline against which the standalone decider must justify its additional cost.}
\label{tab:design}
\centering
\small
\setlength{\tabcolsep}{4pt}
\begin{tabular}{@{}lllll@{}}
\toprule
Decision point & Options & Dataset (test size) & Replaces generation? & Cheap alternative \\
\midrule
Agent control & next tool, or reply & \taubench\ airline (1,076 states; 50 tasks) & No & generator's own call \\
Cache guard & reuse / regenerate & GSM-Plus, LMArena, SearchQueries (2,000 each) & Yes & similarity threshold \\
Router & small / large model & RouterBench (1,998 prompts) & Yes & trained classifier \\
\bottomrule
\end{tabular}
\end{table*}

